% A self-contained theorem/proof fragment for the October 8 continuation.
% Uses only ordinary LaTeX mathematical notation and amssymb's \mathbb.
\subsection{The full scalar commutant and its inherited corners}

Work over $k=\mathbb F_2$. Partition the objects of a finite open complex $C$
by their matching, homological degree, and recovered relative quantum shift.
An allowed scalar change of basis is block diagonal for this partition.
Write $b_1,\ldots,b_t$ for the block sizes and
\[
 N=\sum_a b_a,\qquad V=\sum_a b_a^2,\qquad
 A=\operatorname{End}_{\mathrm{sc}}(C)
   \subseteq\bigoplus_a M_{b_a}(k).
\]
The subscript is essential: these are coefficientwise commuting scalar
endomorphisms. They need not be all endomorphisms in the cobordism category.
Every morphism coefficient contributes its own homogeneous linear
commutation equations. The maintained solver computes a complete basis of
$A$ whenever its configured variable cap permits a computation.

\paragraph{Corner theorem.}
Let $S$ be an allowed invertible change of basis, and suppose the transformed
complex has an actual direct-sum decomposition
$S^{-1}CS=C_1\oplus\cdots\oplus C_s$, with each $C_i$ a whole connected
component of its differential graph. Scalar blocks on a child carry the
inherited parent labels; fresh relative quantum-shift recovery gives the
same blocks on a connected child. Let $e_i$ be the corresponding
coordinate projection. Then
\[
 \operatorname{End}_{\mathrm{sc}}(C_i)
    \cong e_i(S^{-1}AS)e_i.
\]
In particular, restricting the transported members of any complete vector
space basis of $A$ spans the full scalar endomorphism algebra of $C_i$.
This assertion does not require $e_i$ to be central.

\emph{Proof.}
The coordinate inclusion $j_i$ and projection $p_i$ commute with the
transformed differential, and $p_i j_i=1_{C_i}$. Restriction sends a
parent endomorphism $a$ to $p_iS^{-1}aSj_i$. Conversely, if
$u\in\operatorname{End}_{\mathrm{sc}}(C_i)$, then $j_iup_i$, extended by
zero on all other summands, is a typed scalar endomorphism of the transformed
parent. Thus the restriction map is a surjective linear map, and its image
is precisely the displayed corner algebra. Conjugating back by $S$ proves
the assertion about a basis. All changes preserve the parent quantum
labels. On each connected child the recovered quantum labels differ from
these inherited labels by a constant, so the permitted blocks agree.
\hfill$\square$

\paragraph{A necessary algebraic distinction.}
The linear map $a\mapsto e_i a e_i$ from the whole parent algebra to its
corner is not in general an algebra homomorphism. For example, in
$M_2(k)$, with $e=E_{11}$, $a=E_{12}$, and $b=E_{21}$, one has
$eabe=e$ while $eae=ebe=0$. Multiplication on the corner itself is
unambiguous; the algorithm needs the surjective linear map, not a quotient
homomorphism assertion about the entire parent algebra.

\paragraph{Canonical-basis compatibility.}
The existing nullspace solver eliminates equations with their most
significant nonzero coordinates as pivots and lists its free variables in
increasing order. Its output is exactly the unique reduced basis of the
solution space whose pivots are least significant nonzero coordinates.
Therefore reducing the restricted spanning vectors to this latter form
produces the same ordered basis as a fresh child commutation solve.

\emph{Proof.}
Each nullspace basis vector has one free coordinate equal to one and all
other free coordinates zero. The equation pivots determined by that free
coordinate occur only at larger indices. Hence the free coordinate is its
least significant nonzero coordinate, and no such pivot coordinate occurs
in another basis vector. These properties characterize reduced row-echelon
form for the reversed choice of pivot direction and determine the basis
uniquely.\hfill$\square$

This compatibility matters operationally: the bounded basis-prefix search
and the fixed-seed random combinations then inspect exactly the same
candidates, in the same order. General corner inheritance changes the
construction of the search space, not the search policy. With exact
arithmetic and absent an intervening resource interruption, it preserves
the accepted split witnesses and the resulting recursive component order.

\paragraph{Computational scope.}
On a cold recursion tree, general inheritance computes the full scalar
commutant once at the root and derives every subsequently examined child
space by linear algebra, without constructing another commutator system.
An accepted split loaded from an older cache entry need not have a retained
commutant basis; ordinary inheritance can then resume at a descendant.
The cache retains the previous witness format rather than storing large
algebra bases indefinitely.

The implementation has two exact transport kernels. For a low-dimensional
algebra it directly forms $S^{-1}aS$ for its few basis matrices. For a larger
algebra it constructs one packed linear map on matrix variables. The image
of a matrix unit is obtained from
\[
 S^{-1}E_{uv}S=(S^{-1}e_u)(e_v^{\mathsf T}S),
\]
then restricted to all selected children. One application of this packed
linear map transports each parent basis vector. Both kernels give identical
canonical child bases. Their selection affects performance only.

For reference, direct dense transport of a $D$-dimensional algebra costs
$O(D\sum_a b_a^3)$ elementary field operations before child basis
reduction. If child $i$ has $V_i$ scalar variables, reducing $D$ projected
generators costs at most $O(DV_i\min(D,V_i))$ elementary field operations;
packed binary arithmetic improves the corresponding word-operation bound.
These bounds do not imply that transport always beats resolving a sparse
child commutator. The saved quantity is the number of commutator-system
constructions and solves, while the replacement algebra computations have
their own cost. Moreover, cache keys, grading recovery, accepted-witness
verification, and transformation of the differential still inspect
morphism data; the whole recursive scanner is not coefficient independent.

\subsection{A conclusive scalar-locality certificate from dimension saturation}

\paragraph{Saturation theorem.}
Let $A$ be the complete scalar endomorphism algebra, of dimension $D$, and
let $Q\in A$ have actual minimal polynomial $m_Q(t)$. If
\[
 \deg m_Q=D,
\]
then $A=k[Q]\cong k[t]/(m_Q)$. If in addition
\[
 \dim_k\{p\in k[t]/(m_Q):p^2=p\}=1,
\]
then $C$ has no nontrivial scalar direct-sum decomposition.

\emph{Proof.}
The elements $1,Q,\ldots,Q^{D-1}$ are independent by the definition of the
minimal polynomial. They lie in the $D$-dimensional space $A$ and hence form
a basis. Evaluation therefore identifies the whole algebra $A$ with the
displayed polynomial quotient. In characteristic two, squaring is linear
on this commutative quotient. A one-dimensional fixed space contains exactly
$0$ and $1$, so these are its only idempotents. A nontrivial scalar direct
sum would supply a nonzero nonidentity scalar coordinate idempotent, a
contradiction.\hfill$\square$

Equivalently, $m_Q$ has one distinct irreducible factor and the polynomial
quotient is local, including when that factor is repeated. A certificate
does not require that the algebra be a field. The tests include both
$\mathbb F_8$ and $\mathbb F_2[t]/((t^2+t+1)^2)$, exhaustively checking that
their displayed commutants have exactly two idempotents.

Both hypotheses are indispensable operationally. The dimension must be that
of the complete commutant, and the polynomial must be the actual minimal
polynomial, not merely an annihilating polynomial. A primary candidate in a
proper subalgebra proves nothing about idempotents elsewhere in $A$.
For example, the identity of a decomposable complex has primary minimal
polynomial $t+1$ while its full commutant has dimension greater than one.
An explicit negative control verifies that this case is rejected.

The implementation reuses the earlier polynomial-projector routine's exact
minimal polynomial and Berlekamp fixed-space computation. When a projector
search fails and its minimal degree equals the full commutant dimension,
the optional locality flag terminates the otherwise bounded candidate
search with a certificate. This shortcut cannot suppress a possible scalar
split: the certificate establishes that none exists. It makes no assertion
about non-scalar cobordism endomorphisms.

The separate verifier recomputes the entire commutant, verifies commutation
of the proposed generator, recomputes its minimal polynomial, and checks
the fixed-space dimension. It deliberately incurs the full cost again;
certificate auditing is outside benchmark timing. Recorded certificates
include all differential rows, matchings, frontier points, and degrees,
and survive JSON round trips.

\subsection{One-generator inheritance on saturated branches}

\paragraph{Monogenic corner corollary.}
Suppose $A=k[Q]$ has been certified by dimension saturation. For any actual
child coordinate idempotent $e_i$ after an allowed change of basis, its full
endomorphism algebra is generated by the restricted $Q_i$ and its own
identity:
\[
 \operatorname{End}_{\mathrm{sc}}(C_i)=k[Q_i].
\]
Consequently its dimension is $\deg m_{Q_i}$; if $m_{Q_i}$ is primary, the
child is scalar indecomposable without any new commutator solve.

\emph{Proof.}
Conjugation preserves the polynomial-algebra identity. The parent algebra
is commutative, so its child selectors are central. Its corner consists of
the elements $e_ip(Q)e_i=p(Q_i)$, with constants interpreted as multiples
of the corner identity. The corner theorem identifies this with the full
child algebra. Minimal-polynomial dimension and the saturation theorem
now apply.\hfill$\square$

Thus only one generator needs to be transported, rather than a full parent
basis. The implemented optional path recognizes saturation on an accepted
polynomial-projector witness, transports its original generator, and checks
each examined child's minimal polynomial. A local child is finished
immediately. For a nonlocal child, the powers of the restricted generator
are reduced to the same canonical commutant basis as an independent solve,
and the existing bounded splitter resumes.

There is a stronger algorithmic consequence, which should be distinguished
from the implementation just described. Factoring $m_Q$ into its primary
factors, or recursively using the available polynomial projectors, yields a
complete scalar direct-sum decomposition of a saturated monogenic algebra
when split budgets permit. The current implementation completely finishes
the measured two-field branches; it does not replace every nonlocal child's
bounded search by an unlimited recursive factor-decomposition procedure.

\subsection{Controlled families and precise performance claims}

For the existing attached-copy family with $r\geq2$ copies, the commutant is
$M_r(k)$, there are $3r^2$ scalar variables, and each size-$r$ solve builds
$2r^2$ nonzero coefficient equations. With sufficiently large split budget,
the original deterministic policy splits off one copy at each step. Its
total equation construction count is therefore
\[
 2\sum_{j=2}^r j^2
   =\frac{r(r+1)(2r+1)}3-2,
\]
whereas inheritance builds exactly $2r^2$ such equations at the root. This
is an exact cubic-to-quadratic reduction in the number of constructed
commutator equations, not a claim of a corresponding total runtime bound.
The basis candidates and accepted witnesses remain identical. Both the
formula and the witness identity are regression tested through $r=16$.

The separately inherited two-field fixtures have a homogeneous two-term
differential $d=xI+yB$ over a two-dot algebra, where $B$ is conjugate to
companion matrices for two distinct irreducibles of degree $r$. Comparing
the independent $x$ and $y$ coefficients shows that a scalar endomorphism
is a common matrix commuting with $B$. Thus
\[
 \operatorname{End}_{\mathrm{sc}}(C)
  \cong k[t]/(fg)\cong\mathbb F_{2^r}\times\mathbb F_{2^r}.
\]
These are legitimate graded algebraic complexes, without a claimed
realization as knot-scan prefixes. Their root has $8r^2$ scalar variables.
For the recorded dense fixtures at $r=8,10,12$, the accepted generator
saturates the algebra. One root commutant solve, two candidate tests, and
two inherited scalar-locality certificates complete the decomposition.
The earlier primary policy used respectively $46,52,58$ candidates over
the whole recursive compression and three commutant computations.

Natural-knot scans, attached-copy tests, and field-algebra tests have separate
timing scopes. The benchmark retains randomized paired samples and A/A
controls. It compares general inheritance alone, local stopping alone,
their ordinary combination, and one-generator inheritance independently.
No controlled-family speedup establishes a worst-case bound for arbitrary
knot diagrams. In particular, nothing here bounds the number of objects,
matching types, algebra variables, or scalar-indecomposable summands of a
general scan by a quasipolynomial function of its crossing number.

\paragraph{Provenance and proposed follow-up questions.}
The polynomial-projector module and the two-field fixtures are reproduced
unchanged from the earlier October 8 projectors package, with their hashes
recorded separately. Their algorithm is prior work, not a new contribution
of this continuation. The additions are complete-corner transport,
canonical basis preservation, dimension-saturation stopping and verification,
and one-generator inheritance on certified saturated branches.

The next questions are whether to replace the residual bounded scalar search
by constructive finite-algebra radical and semisimple-quotient algorithms;
whether an exact cost model can select transport versus resolving a sparse
commutator; whether additional natural knot families produce saturated or
otherwise tractable scalar algebras; and which structural bounds control
the number and size of scalar-indecomposable components. Extending reuse
across the addition of a crossing is a different problem: a child corner
of the same complex inherits all endomorphisms by extension by zero, while
a newly glued complex can gain endomorphisms that have no ancestor.
